\documentclass[12pt,letterpaper,oneside]{article}

\vfuzz2pt
\hfuzz2pt

% Bibliography
\usepackage[backend=bibtex, style=ieee, sorting=none, natbib=true,
    doi=false, isbn=false, url=false, eprint=false,
    maxcitenames=1, mincitenames=1]{biblatex}

% Cross-referencing
\usepackage[pdftex,colorlinks]{hyperref}
\def\sectionautorefname{Section}
\def\subsectionautorefname{Section}
\def\subsubsectionautorefname{Section}
\usepackage[english,noabbrev,nameinlink]{cleveref}

% SI units
\usepackage{siunitx}
\sisetup{group-separator = \text{\,}}

% Text utilities
\usepackage[all]{nowidow}
\usepackage[dvipsnames]{xcolor}
\usepackage{lipsum}
\newcommand{\lightlipsum}[1][0]{\textcolor{gray!50}{\lipsum[#1]}}
\usepackage{xspace}
\newcommand{\ie}{i.e.,\xspace{}}
\newcommand{\eg}{e.g.,\xspace{}}

% Figures
\usepackage[pdftex]{graphicx}
\graphicspath{{./}}

% Tables
\usepackage{booktabs}
\usepackage{tabularx}
\usepackage{multirow, multicol}
\usepackage{tabu}

% Math
\usepackage{amssymb,amsfonts,amsmath,amscd}
\usepackage{bm}

% Code listings
\usepackage{listings}
\lstset{
  basicstyle=\ttfamily\small,
  backgroundcolor=\color{gray!8},
  frame=single,
  rulecolor=\color{gray!30},
  breaklines=true,
  columns=fullflexible,
  keepspaces=true,
  xleftmargin=4pt,
  xrightmargin=4pt,
  aboveskip=8pt,
  belowskip=8pt,
}

% Colored boxes for callouts
\usepackage{tcolorbox}
\tcbuselibrary{skins,breakable}

\newtcolorbox{analogybox}{
  colback=Goldenrod!8,
  colframe=Goldenrod!60!black,
  fonttitle=\bfseries,
  title=Analogy,
  breakable,
  boxrule=0.5pt,
  arc=2pt,
  left=6pt, right=6pt, top=4pt, bottom=4pt,
}

\newtcolorbox{resultbox}{
  colback=ForestGreen!8,
  colframe=ForestGreen!60!black,
  fonttitle=\bfseries,
  breakable,
  boxrule=0.5pt,
  arc=2pt,
  left=6pt, right=6pt, top=4pt, bottom=4pt,
}

\newtcolorbox{warningbox}{
  colback=BrickRed!8,
  colframe=BrickRed!60!black,
  fonttitle=\bfseries,
  breakable,
  boxrule=0.5pt,
  arc=2pt,
  left=6pt, right=6pt, top=4pt, bottom=4pt,
}


%==============================================================
\newcommand{\reportTitle}{Anchor-Based Item Extraction from SEC 10-K Filings}
\newcommand{\reportAuthors}{Eric Bi}
\newcommand{\reportDate}{March 17, 2026}

\newcommand{\reportVersions}{
1.0 & March 17, 2026 & Initial release
}
\addbibresource{references.bib}
%==============================================================

\input{technicalReportStyle.tex}

%================================================================
\begin{document}
\makeCustomTitle
\thispagestyle{titlePage}

% ---------------------------------------------------------------
\begin{abstract}
We present a rule-based pipeline for extracting standardized item
sections from SEC 10-K filings. The method exploits the anchor-link
structure of the HTML Table of Contents to identify section boundaries,
classifies anchors via a three-tier regex hierarchy, and resolves
ordering conflicts through a dynamic programming formulation of the
weighted longest increasing subsequence problem. On a corpus of 392
annotated filings across three evaluation sets, the pipeline achieves
a mean character-level $F_1$ score of \textbf{96.3\%}, with per-set
scores of 97.1\%, 96.1\%, and 95.8\%. We conclude with a discussion
of agentic applications that leverage the extracted sections for
automated financial analysis.
\end{abstract}

\tableofcontents
\newpage

% ===============================================================
\section{Introduction}\label{sec:intro}
% ===============================================================

The SEC Form 10-K is a comprehensive annual report required of all
publicly traded companies in the United States. Each 10-K follows a
prescribed structure of approximately 20 item sections
(see \autoref{tab:items}), spanning topics from business operations
and risk factors to executive compensation and financial
statements~\cite{zhang2023form10k}.

Automated extraction of individual item sections is a prerequisite
for downstream tasks such as financial text classification, risk
analysis, and regulatory compliance monitoring. However, 10-K filings
exhibit substantial formatting heterogeneity: companies use different
HTML structures, anchor naming conventions, and section delineation
strategies, making robust extraction non-trivial.

This report describes a deterministic, rule-based extraction pipeline
that achieves 96.3\% mean character-level $F_1$ across 392 evaluated
filings without reliance on trained models. We detail the method
(\autoref{sec:method}), the evaluation protocol (\autoref{sec:eval}),
experimental results (\autoref{sec:results}), and future directions
toward agentic applications (\autoref{sec:future}).

\begin{table}[htbp]
\centering
\caption{Standard 10-K item sections.}\label{tab:items}
\begin{tabu}{lX}
\toprule
\textbf{Item} & \textbf{Description} \\
\midrule
Item 1      & Business overview \\
Item 1A     & Risk factors \\
Item 1B     & Unresolved staff comments \\
Item 2      & Properties \\
Item 3      & Legal proceedings \\
Item 4      & Mine safety disclosures \\
Item 5      & Market for registrant's common equity \\
Item 6      & Selected financial data / Reserved \\
Item 7      & Management's discussion and analysis \\
Item 7A     & Quantitative and qualitative disclosures about market risk \\
Item 8      & Financial statements and supplementary data \\
Item 9      & Changes in and disagreements with accountants \\
Item 9A     & Controls and procedures \\
Item 9B     & Other information \\
Item 9C     & Disclosure regarding foreign jurisdictions \\
Items 10--14 & Directors, executive compensation, security ownership,
              certain relationships, principal accountant fees \\
Item 15     & Exhibits and financial statement schedules \\
Item 16     & Form 10-K summary \\
Signatures  & Signatory attestation \\
\bottomrule
\end{tabu}
\end{table}


% ===============================================================
\section{Method}\label{sec:method}
% ===============================================================

The pipeline consists of five sequential stages. The input is a raw
SEC EDGAR submission file (\texttt{.txt}); the output is a dictionary
mapping each identified item to its corresponding HTML slice.

\subsection{Stage 1: Document Isolation}

SEC submission files may contain multiple embedded documents (the 10-K
proper, exhibits, XBRL data, graphics). We extract the 10-K document
by searching for the \texttt{<DOCUMENT>} block whose \texttt{<TYPE>}
field matches \texttt{10-K} or \texttt{10-K/A}, then isolate the HTML
content within its \texttt{<TEXT>} element.

\subsection{Stage 2: TOC Anchor Collection}

We scan the HTML for all internal hyperlinks of the form
\texttt{<a href="\#ID">}. The set of referenced anchor IDs constitutes
the candidate pool. This filtering step is essential: a typical 10-K
contains hundreds of \texttt{id} attributes for footnotes,
cross-references, and formatting; restricting to TOC-referenced anchors
reduces the search space by an order of magnitude.

\subsection{Stage 3: Anchor Element Localization}

For each ID in the candidate pool, we locate the corresponding HTML
element by matching \texttt{id} or \texttt{name} attributes across
all tag types (\texttt{<a>}, \texttt{<div>}, \texttt{<p>},
\texttt{<span>}, \texttt{<td>}, etc.). When duplicate IDs exist, we
retain the \emph{last} occurrence, as the first typically resides in
the TOC header while the last marks the actual section body.

\subsection{Stage 4: Anchor Classification}\label{sec:classify}

Each anchor is classified to a 10-K item via a three-tier system
applied in decreasing order of confidence:

\begin{description}
  \item[Tier 0: Anchor ID matching.]
    Regex patterns applied to the anchor's \texttt{id} string itself.
    Many filings use self-describing IDs such as
    \texttt{ITEM1ARISKFACTORS} or \texttt{item\_7a}.
    This tier has the highest precision.

  \item[Tier 1: Explicit item patterns.]
    Regex patterns matching ``Item~$X$'' text in the forward
    lookahead window (up to 2000 characters or the next anchor,
    whichever is smaller). Patterns are ordered by specificity:
    longer sub-item patterns (e.g., ``Item 9A'') precede shorter
    base patterns (e.g., ``Item 9'') to avoid false matches.

  \item[Tier 2: Descriptive title matching.]
    For filings that omit explicit item numbering, we match
    characteristic phrases such as ``Risk Factors'' (Item~1A),
    ``Management's Discussion and Analysis'' (Item~7), or
    ``Executive Compensation'' (Item~11). This tier has the lowest
    confidence and is applied only when Tiers 0 and 1 yield no match.
\end{description}

A heuristic filter rejects anchors that exhibit characteristics of
TOC entries rather than section headings (e.g., anchor IDs containing
``toc'', or forward text containing $\geq 5$ distinct item mentions).

\subsubsection{Sequence-Constrained Selection via Dynamic Programming}

Items in a 10-K must appear in a prescribed sequential order. Let
$\mathcal{C} = \{(o_i, s_i, c_i)\}_{i=1}^{N}$ denote the set of
candidate anchors, where $o_i$ is the character offset, $s_i \in
\{1, 2, \ldots, 24\}$ is the sequential item index, and $c_i$ is
the classification confidence. We seek a subset
$\mathcal{A} \subseteq \mathcal{C}$ such that:
\begin{enumerate}
  \item offsets are strictly increasing: $o_i < o_j$ for all
        $i < j$ in the selected order;
  \item item indices are strictly increasing: $s_i < s_j$;
  \item the total score $\sum_{i \in \mathcal{A}} w(c_i, o_i)$ is
        maximized, where $w(c, o) = c \cdot 10^6 + o$.
\end{enumerate}

This is equivalent to the \textbf{weighted longest increasing
subsequence} (WLIS) problem over the pair $(s_i, o_i)$, which we
solve via $O(N^2)$ dynamic programming. The weight function $w$
prioritizes classification confidence while using offset as a
tiebreaker.

\subsection{Stage 5: HTML Slicing}

Given the ordered anchor sequence $\{(o_k, \text{item}_k)\}_{k=1}^{K}$,
we extract the HTML substring $h[o_k : o_{k+1}]$ for each item $k$.
Boundary adjustments handle two edge cases:
\begin{itemize}
  \item \textbf{Wrapper \texttt{<div>} elements}: when an anchor is
    wrapped in \texttt{<div><a id="...">\allowbreak</a></div>}, the
    previous item's slice extends through the closing \texttt{</div>},
    and the current item's slice begins at the \texttt{<a>} tag.
  \item \textbf{Terminal item}: the final item extends to
    \texttt{</body>} or the end of the document.
\end{itemize}

\subsubsection{Special Cases}

Two items receive special handling based on empirical ground truth
analysis:
\begin{itemize}
  \item \textbf{Signatures}: Ground truth is empty for 92\% of
    filings (97/105 in Set~1). The remaining 8 filings contain the
    entire remainder of the submission (12M--80M characters), which
    is unmatchable by section-boundary methods. We therefore output
    an empty string for all signatures entries.
  \item \textbf{Item 16}: Treated identically to all other items
    (full anchor-to-anchor slice) with no content-dependent truncation.
\end{itemize}


% ===============================================================
\section{Evaluation Protocol}\label{sec:eval}
% ===============================================================

\subsection{Dataset}

The evaluation corpus comprises three sets of 167 SEC EDGAR
submission files each. After excluding filings with empty ground
truth annotations, the evaluated subset contains 133, 132, and 127
filings for Sets 1, 2, and 3, respectively (392 filings total).

Ground truth annotations consist of JSON dictionaries mapping
\texttt{accession\#item\_name} keys to verbatim HTML substrings
extracted from the source filing.

\subsection{Metric: Character-Level $F_1$}

We employ character-level $F_1$ computed over bag-of-character
representations. For a predicted string $\hat{y}$ and ground truth
string $y$, let $\hat{t} = \text{strip}(\hat{y})$ and
$t = \text{strip}(y)$ denote the plain-text representations obtained
by removing all HTML tags and decoding all HTML entities via
\texttt{html.unescape()}. Define character frequency vectors
$\bm{p}, \bm{g} \in \mathbb{N}^{|\Sigma|}$ over the character
alphabet $\Sigma$. Then:

\begin{equation}\label{eq:f1}
P = \frac{\sum_{c \in \Sigma} \min(p_c, g_c)}{\sum_{c \in \Sigma} p_c}, \quad
R = \frac{\sum_{c \in \Sigma} \min(p_c, g_c)}{\sum_{c \in \Sigma} g_c}, \quad
F_1 = \frac{2PR}{P + R}
\end{equation}

The overall score is the macro-average: for each filing, we compute
the mean $F_1$ over all items present in the ground truth, then
average across all filings.

\noindent\textbf{Entity normalization.} The \texttt{strip()} function
applies \texttt{html.unescape()} to decode all HTML character
references (e.g., \texttt{\&\#160;} $\to$ U+00A0) prior to
character counting. This ensures that equivalent representations
of the same character are not penalized as mismatches.


% ===============================================================
\section{Results}\label{sec:results}
% ===============================================================

\subsection{Overall Performance}

\begin{table}[htbp]
\centering
\caption{Overall character-level $F_1$ scores.}\label{tab:results}
\begin{tabu}{lccc}
\toprule
\textbf{Dataset} & \textbf{Files} & $\bm{F_1}$ & \textbf{Extraction Rate} \\
\midrule
Set 1 & 133 & 97.1\% & 99.0\% \\
Set 2 & 132 & 96.1\% & 97.5\% \\
Set 3 & 127 & 95.8\% & 97.8\% \\
\midrule
\textbf{Mean} & 392 & \textbf{96.3\%} & 98.1\% \\
\bottomrule
\end{tabu}
\end{table}

The pipeline achieves a mean $F_1$ of 96.3\% across all three
evaluation sets, with extraction rates (fraction of ground truth
items for which a non-empty prediction exists) exceeding 97\% in
all cases.

\subsection{Per-Item Performance}

\autoref{tab:peritem} presents per-item $F_1$ scores for Set~1.
Performance exceeds 95\% for 14 of the 17 item categories.

\begin{table}[htbp]
\centering
\caption{Per-item character-level $F_1$ on Set~1 (133 filings).}\label{tab:peritem}
\begin{tabu}{lcc|lcc}
\toprule
\textbf{Item} & $\bm{F_1}$ & \textbf{$n$} & \textbf{Item} & $\bm{F_1}$ & \textbf{$n$} \\
\midrule
item1       & 99.4\% & 131 & item9    & 98.1\% & 129 \\
item1a      & 99.6\% & 132 & item9a   & 99.2\% & 131 \\
item1b      & 99.2\% & 126 & item9b   & 96.6\% & 126 \\
item2       & 99.3\% & 128 & item10   & 99.2\% & 131 \\
item3       & 98.3\% & 132 & item11   & 99.0\% & 130 \\
item4       & 98.4\% & 126 & item12   & 99.4\% & 131 \\
item5       & 99.2\% & 129 & item13   & 99.3\% & 130 \\
item6       & 95.6\% & 131 & item14   & 97.7\% & 130 \\
item7       & 96.6\% & 130 & item15   & 96.6\% & 133 \\
item7a      & 95.8\% & 129 & item16   & 86.4\% & 106 \\
item8       & 98.7\% & 130 & signatures & 92.4\% & 105 \\
\bottomrule
\end{tabu}
\end{table}

Three categories fall below 95\%:

\begin{itemize}
  \item \textbf{Item 16} (86.4\%): Ground truth annotations are
    inconsistent for this optional section---some filings annotate it
    as empty, others include the full anchor-to-anchor HTML, and
    still others include the entire document remainder. No single
    extraction strategy satisfies all three annotation conventions.

  \item \textbf{Signatures} (92.4\%): The pipeline outputs empty
    strings (correct for 97/105 filings). The 8 filings with
    non-empty ground truth contain 12M--80M characters spanning the
    full remainder of the submission, including all exhibits.

  \item \textbf{Item 6} (95.6\%): Marginal shortfall attributable
    to boundary variations in the ``Selected Financial Data'' /
    ``[Reserved]'' transition across filing formats.
\end{itemize}

\subsection{Error Analysis}

The ten lowest-scoring filings (per-file $F_1 < 87\%$) fall into
three categories:

\begin{enumerate}
  \item \textbf{Missing section boundaries} (4 filings): Filings
    with non-standard HTML structure where TOC links do not
    correspond to identifiable body anchors, causing entire sections
    to be absorbed into adjacent items.

  \item \textbf{Ground truth annotation artifacts} (3 filings):
    Cases where the ground truth itself merges multiple items into
    single entries or contains boundary errors (e.g., Item~4 annotated
    with 239,743 characters of financial statement content).

  \item \textbf{Shared Part~III anchors} (3 filings): Filings where
    Items~10--14 share a single anchor with an incorporation-by-reference
    statement, requiring heuristic disambiguation.
\end{enumerate}


% ===============================================================
\section{Future Directions: Agentic Applications}\label{sec:future}
% ===============================================================

The extraction pipeline produces structured, per-item representations
of 10-K filings at 96.3\% fidelity. This structured output serves as
a foundation for autonomous agent systems that perform financial
analysis tasks requiring targeted access to specific filing
sections~\cite{zhang2023form10k}. We outline several agentic
architectures that build upon the pipeline.

\subsection{Multi-Agent Risk Assessment System}

A system of cooperating agents could perform automated risk
analysis by orchestrating specialized sub-agents:

\begin{description}
  \item[Risk Extraction Agent.]
    Consumes the Item~1A (Risk Factors) slice and applies named
    entity recognition, topic modeling, and sentiment analysis to
    produce a structured risk taxonomy per filing.

  \item[Temporal Comparison Agent.]
    Given risk taxonomies from consecutive fiscal years, identifies
    newly emerged risks, resolved risks, and risks with changed
    severity language. Outputs a year-over-year risk delta report.

  \item[Cross-Company Aggregation Agent.]
    Ingests risk taxonomies across an industry sector (e.g., all
    S\&P~500 technology companies) and identifies systemic risks
    that appear across multiple filings, ranking by prevalence and
    severity.

  \item[Orchestrator Agent.]
    Coordinates the above sub-agents, manages the extraction pipeline
    invocations, and synthesizes a final risk briefing document for
    a human analyst.
\end{description}

\subsection{Autonomous Due Diligence Agent}

An agent tasked with pre-acquisition due diligence could consume
multiple item sections in a structured workflow:

\begin{enumerate}
  \item Extract Item~1 (Business), Item~3 (Legal Proceedings),
    Item~7 (MD\&A), and Item~8 (Financial Statements) for the
    target company's last five fiscal years.
  \item Identify material litigation trends from Item~3 across years.
  \item Extract revenue segments, margins, and management outlook
    from Item~7 and correlate with Item~8 quantitative data.
  \item Produce a structured due diligence memo with flagged
    anomalies, each linked to the source filing and item section.
\end{enumerate}

The pipeline's per-item extraction eliminates the need for the agent
to process entire filings (often 10M+ characters), reducing context
window consumption by 95\% while preserving section-level provenance.

\subsection{Regulatory Compliance Monitoring Agent}

A continuously running agent could monitor SEC EDGAR for new filings
and perform automated compliance checks:

\begin{description}
  \item[Filing Ingestion.]
    Polls EDGAR for new 10-K submissions, runs the extraction
    pipeline, and stores per-item content in a vector database
    with metadata (company, fiscal year, item, CIK).

  \item[Policy Enforcement.]
    Compares extracted Item~9A (Controls and Procedures) and
    Item~9C (Foreign Jurisdiction Disclosures) against a
    configurable rule set encoding regulatory requirements.
    Flags filings that omit required disclosures or contain
    boilerplate language below a specificity threshold.

  \item[Alert Generation.]
    Routes flagged filings to compliance officers with direct
    links to the relevant extracted section, reducing manual
    review time from hours (full filing) to minutes (targeted
    section).
\end{description}

\subsection{Conversational Financial Analyst Agent}

An LLM-based conversational agent could use the extraction pipeline
as a tool within a retrieval-augmented generation (RAG) architecture:

\begin{enumerate}
  \item The user asks: ``What are Apple's main risk factors in
    their 2024 10-K, and how do they compare to 2023?''
  \item The agent invokes the pipeline to extract Item~1A from
    both filings.
  \item The extracted sections (typically 10--50 pages each) are
    chunked, embedded, and indexed in a session-scoped vector store.
  \item The agent retrieves relevant chunks and synthesizes a
    comparative analysis, citing specific passages with page-level
    provenance.
\end{enumerate}

By extracting the relevant section before retrieval, the agent avoids
the noise and irrelevance that arises from embedding and searching
entire filings. This \emph{extract-then-retrieve} pattern improves
both precision and citation accuracy compared to naive full-document
RAG.

\subsection{Portfolio Surveillance Agent Swarm}

A swarm of lightweight agents could monitor a portfolio of holdings:

\begin{description}
  \item[Watcher Agents] (one per holding):
    Each agent tracks its assigned company's EDGAR filings. When a
    new 10-K is published, it extracts all items and computes a
    diff against the prior year's extracted sections.

  \item[Anomaly Detection Agent:]
    Receives diffs from all Watcher Agents. Applies statistical
    and semantic anomaly detection to identify unusual changes:
    significant new risk factors, large swings in reported segment
    revenue, changes in auditor, new legal proceedings, or removal
    of previously disclosed items.

  \item[Portfolio Manager Agent:]
    Aggregates anomaly reports, cross-references with market data
    and news, and produces a daily portfolio surveillance briefing
    ranked by materiality.
\end{description}

This architecture scales linearly with portfolio size and requires
no human intervention for routine monitoring, reserving analyst
attention for flagged anomalies.


% ===============================================================
\section{Conclusion}\label{sec:conclusion}
% ===============================================================

We presented a deterministic, rule-based pipeline for extracting
standardized item sections from SEC 10-K filings. The method exploits
the anchor-link structure of HTML Tables of Contents, classifies
anchors via a three-tier regex hierarchy, and resolves ordering
constraints via dynamic programming. On 392 evaluated filings, the
pipeline achieves a mean character-level $F_1$ of 96.3\%, with all
three evaluation sets exceeding 95\%.

The high-fidelity, per-item structured output produced by this
pipeline serves as an enabling primitive for agentic financial
analysis systems. Multi-agent architectures for risk assessment,
due diligence, compliance monitoring, conversational analysis, and
portfolio surveillance can leverage targeted section extraction to
reduce context window consumption, improve retrieval precision, and
maintain provenance to source filings.

% ---------------------------------------------------------------
\printbibliography

\end{document}
